\scanpage{254}

% folio and opening word
\mspage{7}donc

% opening comparison
\[f_*(\mathrm{hom}_{\mathbb D}(\unsure{\xi|\zeta},F))\simeq q_*q^*(f_*(F))\overset{\mathrm{def}}{=}\mathrm{hom}_{\mathbb D}(\xi,f_*(F))\]

% below opening comparison
qed. L'énoncé pour $F\mapsto\xi\times_{\mathbb D}F$ est dual de celui qu'on vient de prouver.

% upper middle, before adjunction
Reprenons, pour $\xi\in(\mathrm{Hot})_0$, la formule d'adjonction

% adjunction, label left
\[\mathrm{Hom}(\xi\times_{\mathbb D}M,N)\simeq\mathrm{Hom}(M,\unsure{\mathrm{Hom}|\mathrm{hom}}_{\mathbb D}(\xi,N))\overset{\mathrm{def}}{=}\mathrm{Hom}_{\xi}(M,N)\tag*{\(\unsure{a|\alpha})\)}\]

% struck beginning immediately below adjunction
\deleted{Pour $M,N$ fixés, et}

% above Hom formula
Si $\xi$ est défini par $X$ dans $\mathrm{Diag}$, les deux \unsure{mbs|mbr} de (a) s'identifient à

% middle, Hom in the fibre
\[\mathrm{Hom}_{\mathbb D(X)}(p^*M,p^*N),\]

% middle, structural functor
où $p:X\longrightarrow e$ est le foncteur structural.

% middle, variable coefficient
Fixons $M,N$, et considérons l'expression (a) comme contrafoncteur

% target of contrafunctor
\[(\mathrm{Hot}_0)^{\unsure{\circ|0}}\longrightarrow\mathrm{Ens}\]

% bottom paragraph and inserted letter
Il semblerait que dans tous les cas connus, ce foncteur soit ``représentable'' par un objet, \unsure{sinon} de $\mathrm{Hot}_0$ lui-\unsure{m\textsuperscript{=}|même}, du moins \deleted{\unsure{un}}objet dans $\mathrm{Hot}$. \deleted{\unsure{Un}} Si \deleted{\unsure{l'|a}}on ne peut le définir comme objet \inserted{$H$} de $\mathrm{Hot}_0$, le fait qu'on ait un isom canonique

% bottom formula (10)
\[\mathrm{Hom}_{\xi}(M,N)\simeq\mathrm{Hom}_{\mathrm{Hot}}(\xi,H),\tag*{(10)}\]

% last line, continuing on 255
fonctoriel en $\xi\in\mathrm{Ob}\,\mathrm{Hot}_0$, ne définit pas \unsure{forcément}

% archive sheet number: 254
